% Round 17. Root. Independent review required before activation.
\section{Ordinary quadratic models without a supplied exact equation}
\label{r17on:section}

QF17 constructs native semistable chains and numerical
admissible extensions from exact pointed etale equations.
Here we verify that the exact equation is automatic at an
ordinary quadratic special-fiber point in odd residue
characteristic. The local normal-form principle is classical;
the purpose of the proof is to retain native residue fields
and to make its use in the all-degree actual theorem precise.
In particular, no formal coordinate change is silently
algebraized.

Throughout this section $R=K^\circ$ is a complete discrete
valuation ring of mixed characteristic $(0,p)$, $p\ne2$,
with uniformizer $\pi$ and finite residue field $k$.
A special-fiber point of relative dimension $n$ is called
\emph{ordinary quadratic} when its completed geometric local
ring is a hypersurface in $n+1$ variables with zero linear
term and nondegenerate quadratic term. Equivalently its
projectivized tangent cone is a smooth quadric. The residue
fields at the closed points are retained rather than split.

\begin{lemma}[A smooth local ambient is automatic]
\label{r17on:ambient}
Let $\mathcal X/R$ be flat of finite type, with a closed
ordinary quadratic point $p$ of its special fiber, of
relative dimension $n$. After passage to a
residue-isomorphic pointed etale neighborhood of $p$,
$\mathcal X$ is a Cartier hypersurface in a smooth
$R$-scheme of relative dimension $n+1$.
One may take this neighborhood over the finite unramified
extension $R_p/R$ with residue field $\kappa(p)$, choosing
the diagonal point after base change; its map to the
original scheme remains pointed etale with residue
field $\kappa(p)$.
\end{lemma}

\begin{proof}
First make the indicated unramified base change and
choose the point with the original residue field. This
is an etale change, not a decision to replace the global
model field. Work over $R_p$ and write this point as a
rational point of the special fiber.

Choose a local affine presentation in a smooth affine
space over $R_p$, of relative dimension $M$. At the
chosen point write its regular local ring as $B$,
its defining ideal as $I$, and $A=B/I$.
Flatness of $A/R_p$ and of $B/R_p$ gives
\[
 I/\pi I\ \simeq\
  \ker(B/\pi B\longrightarrow A/\pi A).
 \tag{ON17.1}\label{r17on:tor}
\]
Indeed the left injection is precisely the vanishing of
$\operatorname{Tor}_1^{R_p}(A,k_p)$.
The special local ring is a complete intersection of
codimension $c=M-n$, and its embedding dimension is
$n+1$. To see the same statements before completion,
use faithful flatness of completion: minimal numbers
of generators and the complete-intersection property
of an ideal in a regular local ring are detected there.
Thus the special ideal has $c$ generators whose linear
parts have rank $c-1$.

Lift these generators to $f_1,\ldots,f_c\in I$.
Equation~\eqref{r17on:tor} and local Nakayama give
$I=(f_1,\ldots,f_c)$. After a constant invertible
linear combination, the first $c-1$ special differentials
are independent. The Jacobian criterion makes
\[
                  W=V(f_1,\ldots,f_{c-1})
\]
smooth of relative dimension $n+1$ on a neighborhood of
the point. The last equation $f=f_c|_W$ is nonzero in
its regular local ring: otherwise the special fiber
would have dimension $n+1$. Therefore it is a
nonzerodivisor, and $\mathcal X=V(f)\subset W$ is
Cartier near the point. Its special quadratic term is
the original nondegenerate one, after the elimination
of the independent linear equations.

The lci lifting facts used here can also be read in
Stacks Project, Lemma 29.31.10 (Tag \texttt{01UE}) and
Section 37.62 (Tag \texttt{068E}). The direct Tor,
Nakayama and Jacobian argument records the relative
codimension-one conclusion needed below.
\end{proof}

\begin{proposition}[Native pointed etale quadratic normal form]
\label{r17on:normal-form}
In the situation of Lemma~\ref{r17on:ambient}, assume
the generic fiber of $\mathcal X$ is smooth.
There are a nondegenerate quadratic form
$Q\in R_p[z_0,\ldots,z_n]$, a unit $u\in R_p^\times$
and an integer $b\ge1$ such that the local model has a
common residue-isomorphic pointed etale neighborhood with
\[
       \operatorname{Spec}
       R_p[z_0,\ldots,z_n]/(Q(z)-u\pi^b)
 \tag{ON17.2}\label{r17on:exact}
\]
at its origin. The residue quadratic form is the
original tangent-cone form up to an invertible change
of coordinates and a unit. No additional extension of
$\kappa(p)$ is required.
\end{proposition}

\begin{proof}
Use the smooth ambient $W/R_p$ and function $f$ from
the preceding lemma, shrinking around the chosen point
whenever necessary. Choose relative etale coordinates
$t_0,\ldots,t_n$. Express
$df=\sum_i F_i\,dt_i$. At the special point all $F_i$
vanish, and the Jacobian of $(F_i)$ is the nondegenerate
Hessian. Therefore the relative critical locus
$V(F_0,\ldots,F_n)$ is etale over $R_p$ at that point.
The henselian property supplies a section $a$ through
it. Its values are native elements of $R_p$.

Replace $t_i$ by $x_i=t_i-t_i(a)$.
The section ideal is generated by the $x_i$ after
shrinking: in an etale coordinate neighborhood the
section is the component of the zero coordinate fiber
through the chosen point. Set $c=f(a)\in R_p$.
Since $df|_a=0$, the conormal sequence for the smooth
section gives
\[
                   f-c\in (x_0,\ldots,x_n)^2.
 \tag{ON17.3}\label{r17on:quadratic-ideal}
\]
The section lies over the special zero of $f$, so
$c\in\pi R_p$. If $c=0$, its generic point would
be a critical point on the generic hypersurface,
contradicting smoothness. Hence $c\ne0$.

Because $2$ is invertible, write the exact identity
\[
                      f-c=x^{\mathsf T}H(x)x
\]
with a symmetric matrix of regular functions.
Its value at the section is a unimodular symmetric
matrix over $R_p$. A constant congruence over this
local ring diagonalizes it with unit entries
$\epsilon_0,\ldots,\epsilon_n$. Such diagonalization
requires no splitting of the quadratic form: reduction
modulo $\pi$ has a vector of nonzero norm, one splits
off that rank-one summand and continues on its
unimodular orthogonal complement, then lifts the
chosen bases.

After that constant change, all leading principal
minors are units along the section. Shrink so they
are invertible functions. The finite Gaussian
congruence (or $LDL^{\mathsf T}$ elimination) over
the ring of regular functions gives
\[
     H(x)=T(x)^{\mathsf T}
          \operatorname{diag}(d_0(x),\ldots,d_n(x))T(x),
 \qquad T(a)=1,\quad d_i(a)=\epsilon_i.
 \tag{ON17.4}\label{r17on:ldl}
\]
All denominators are those unit principal minors.
In particular this is a finite algebraic operation.

Adjoin variables $s_i$ satisfying
$s_i^2=d_i(x)/\epsilon_i$ and choose the branch with
$s_i(a)=1$. This is a finite etale cover because
$2s_i$ is invertible; its selected section has
unchanged residue field. Put
\[
                  z_i=s_i\,(T(x)x)_i.
 \tag{ON17.5}\label{r17on:new-coordinates}
\]
Then, as an equality of regular functions,
\[
                   f-c=\sum_i\epsilon_i z_i^2.
 \tag{ON17.6}\label{r17on:identity}
\]
The relative Jacobian of $z$ along the section is
the identity after the prior constant basis change.
Thus the map $z$ to affine $(n+1)$-space over $R_p$
is etale on a pointed neighborhood.

Writing $Q=\sum_i\epsilon_i z_i^2$ and
$-c=u\pi^b$, its restriction to $f=0$ is the
base change of that etale map to $Q=u\pi^b$.
The same neighborhood is etale over the original
hypersurface. Both maps have the chosen residue
field $\kappa(p)$. This is the claimed common
pointed etale normal form.
No infinite series substitution, analytic inverse
function theorem, or unproved passage from a formal
isomorphism to an etale one has been used.
\end{proof}

\begin{theorem}[Full actual G for ordinary quadratic reduction]
\label{r17on:actual}
Let $X/K$ be smooth projective geometrically integral
of dimension $n\ge3$, where $K/\mathbf Q_p$ is finite
and $p\ne2$. Suppose it has an integral flat projective
model $\mathcal X/R$ whose special fiber is geometrically
integral and smooth except for finitely many ordinary
quadratic points. The total model need not be regular,
and no bound is imposed on the orders of its node
deformations.

Then every geometrically homologically trivial rational
cycle in every positive codimension satisfies strong G
over the original field $K$, on the complement of its
exact specified support and against every original compact
actual test.
The same assertion holds for any original-field smooth
projective target admitting a projective dominant witness
of this reduction type after a finite auxiliary extension.
\end{theorem}

\begin{proof}
Flatness and smoothness of the special fiber away from
the indicated points make the model morphism smooth
there. At each point, Proposition~\ref{r17on:normal-form}
supplies exactly the native pointed etale datum of
Definition~\ref{r17qf:thick-datum}, with an arbitrary
positive integer $b_j$. The native unramified rings
$R_{p_j}$ describe only the pointed local charts;
they are not multiplied together as extra global
constant extensions.

Apply QF17. A single ramified quadratic extension
$K(\sqrt\pi)/K$ followed by the stated reduced-point
blowups gives a regular reduced quadric forest.
Its native numerical correction exists in every
degree, hence its actual obstruction is zero.
The frozen dominant and full finite-constant
detectors descend that zero obstruction with all
degrees intact. The curved pair for the final
specified cycle reconstructs the required ordinary
actual form with its exact support.
The dominant-witness assertion is the same frozen
detector and does not push the form through an
alteration or assume that alteration is finite
away from the cycle.
\end{proof}

\begin{corollary}[Finite extensions preserve this proved reduction sector]
\label{r17on:extensions}
Under the hypotheses of Theorem~\ref{r17on:actual},
every finite extension $L/K$ has the full strong-G
conclusion for all geometrically homologically trivial
rational cycles on $X_L$, including cycles not defined
over $K$.
\end{corollary}

\begin{proof}
Base change preserves the geometric integral special
fiber and its ordinary quadratic points. It preserves
flatness and the smooth geometrically integral generic
fiber, and the base-changed model is integral since
it is flat with integral generic fiber. In the
native local normal form, writing
$\pi_K=v\pi_L^{e(L/K)}$ replaces the node order $b$
by $e(L/K)b$ and multiplies its unit coefficient
by $v^b$. Nondegeneracy of the special quadratic
form is unchanged. Residue extensions are handled
at each native point without assuming splitness.
The theorem therefore applies afresh to the whole
Chow group over $L$. This is stronger than
ascent of cycles already defined over $K$.
\end{proof}

\begin{corollary}[Complete height for the specified ordinary-reduction sector]
\label{r17on:height}
Let $X/F$ be smooth projective geometrically integral
of dimension $n\ge3$ over a number field.
Suppose that every finite completion has, after a finite
auxiliary local extension if necessary, a projective
dominant witness with either good reduction or the
ordinary quadratic reduction of
Theorem~\ref{r17on:actual}. Require odd residue
characteristic for the latter alternative.
Then the frozen height construction gives the canonical
real Beilinson--Bloch pairing on the entire geometric
homological rational Chow groups in complementary degrees:
\[
 \operatorname{CH}^r(X)_{\mathbf Q,\mathrm{geom\,hom}}
 \times
 \operatorname{CH}^{n+1-r}(X)_{\mathbf Q,\mathrm{geom\,hom}}
                         \longrightarrow\mathbf R.
 \tag{ON17.7}\label{r17on:height-pairing}
\]
The pairing has the frozen rational-equivalence,
choice-independence, correspondence-adjunction and
number-field normalization properties. GA17 identifies
it with normalized admissible arithmetic intersection
on any fixed regular arithmetic alteration as there.
\end{corollary}

\begin{proof}
The good-reduction detector and
Theorem~\ref{r17on:actual} give obstruction zero for
all local geometric-homological cycles in every degree.
Thus the global domain $\mathcal H^r(X)$ is the full
geometric-homological Chow group for this specified
class of varieties. Apply the frozen PH16 implication
and the GA17 admissible-alteration comparison.
No additional arbitrary-Green residue premise is
introduced. This is a theorem for the stated
reduction sector, not general G or a general
unconditional height theorem.
\end{proof}

\paragraph{Published antecedent and the exact new use.}
The deformation normal form and its order reduction
are classical: Jannsen--Saito, \emph{Bertini theorems
and Lefschetz pencils over discrete valuation rings},
Section 3, Lemma 4 and Theorem 5, printed pp.11--13,
\url{https://arxiv.org/pdf/0911.1470}.
Their Lemma 4 cites SGA7, Expose XV, 1.3.2.
The proof above keeps an unsplit native residue field
and supplies common pointed etale charts explicitly.
The research conclusion is the activation of all
ordinary quadratic models through QF17's all-degree
native numerical correction and the frozen ordinary
actual detector. No novelty is claimed for the
classical local normal-form statement.
